\chapter{Operador constitutivo bicapa del vacío: positividad, completación \texorpdfstring{\(3\to4\)}{3→4} del modo 30 y reconstrucción espectral de \texorpdfstring{\(\varepsilon\)}{ε}, \texorpdfstring{\(\mu\)}{μ}, \texorpdfstring{\(Z\)}{Z} y \texorpdfstring{\(c\)}{c}.}
\label{chap:p3-83-31-vacio_epsilon_mu_z_c}

\section{Descomposici\'on espectral bicapa del operador angular}

Sea \(R=R^*=R^{-1}\) la involuci\'on activa y

\[
P_\pm=\frac{I\pm R}{2},
\qquad
T=q_+P_+ +q_-P_-,
\qquad 0<q_+<q_-<1.
\]

Con
\[
x=A_{\rm rad}=\frac{\pi A}{180},\qquad
y=C^*_{\rm rad}=\frac{\pi C^*}{180},
\]
los dos canales se construyen antes de evaluar el vac\'io:
\begin{equation}
q_+=e^{-x-y},\qquad q_-=e^{-x+y}.
\label{p3-83-c31-s1:eq:vacuum-angular-channels}
\end{equation}
El intercambio de hojas invierte \(C^*\), es decir \(y\mapsto-y\), y permuta \(q_+\leftrightarrow q_-\). La pareja es as\'i la lectura espectral de un mismo antecedente angular, no dos par\'ametros ajustados por separado.

\begin{definition}[Respuesta constitutiva]
\begin{equation}
\mathcal V_{90,120}=(I-T^{90})(I-T^{120})^{-1}.
\label{p3-83-c31-s1:eq:vacuum-operator}
\end{equation}
\end{definition}

Como \(\|T\|<1\), el denominador es positivo e invertible. El c\'alculo funcional da

\begin{equation}
\mathcal V_{90,120}=r_+P_+ +r_-P_-,
\qquad
r_\pm=\frac{1-q_\pm^{90}}{1-q_\pm^{120}}.
\label{p3-83-c31-s1:eq:vacuum-eigenvalues}
\end{equation}

La afirmaci\'on de invertibilidad merece hacerse expl\'icita. Para cada hoja,
\(0<q_\pm^{120}<1\); por tanto
\[
I-T^{120}=(1-q_+^{120})P_++(1-q_-^{120})P_-
\]
es estrictamente positivo y
\[
(I-T^{120})^{-1}
=\frac{P_+}{1-q_+^{120}}+
 \frac{P_-}{1-q_-^{120}}.
\]
La respuesta \(\mathcal V_{90,120}\) no se obtiene resolviendo cuatro
ecuaciones escalares separadas: es el resultado del c\'alculo funcional de una
sola contracci\'on positiva. Esta observaci\'on impide que permeabilidad,
permitividad, impedancia y propagaci\'on pierdan el antecedente que comparten.

\begin{theorem}[Determinaci\'on espectral conjunta de las cuatro relaciones constitutivas]
Los cuatro caracteres constitutivos
\begin{equation}
\widehat\varepsilon=r_+^2,\qquad
\widehat\mu=r_-^2,\qquad
\widehat Z=\frac{r_-}{r_+},\qquad
\widehat c=\frac1{r_+r_-}
\label{p3-83-c31-s1:eq:vacuum-four}
\end{equation}
son funciones espectrales de \(\mathcal V_{90,120}\). Satisfacen
\begin{equation}
\widehat\mu\widehat\varepsilon=\widehat c^{-2},
\qquad
\frac{\widehat\mu}{\widehat\varepsilon}=\widehat Z^2.
\label{p3-83-c31-s1:eq:vacuum-relations}
\end{equation}
\end{theorem}

\begin{proof}
Las identidades se obtienen sustituyendo \cref{p3-83-c31-s1:eq:vacuum-four}. No interviene ninguna carta dimensional ni ning\'un valor exterior.
\end{proof}

\begin{corollary}[Reconstrucci\'on de las dos hojas]
Los cuatro caracteres no contienen cuatro grados de libertad. Cualquiera de
los pares \((\widehat\mu,\widehat\varepsilon)\),
\((\widehat Z,\widehat c)\) o \((\mathcal E,\mathcal B)\) reconstruye los
dos autovalores positivos:
\begin{align}
r_-&=\sqrt{\widehat\mu}
=\sqrt{\widehat Z/\widehat c},
&r_+&=\sqrt{\widehat\varepsilon}
=\frac1{\sqrt{\widehat Z\widehat c}},
\label{p3-83-c31-s1:eq:vacuum-reconstruct-r}\\
\log r_-&=\frac{\mathcal E+\mathcal B}{2},
&\log r_+&=\frac{\mathcal E-\mathcal B}{2}.
\label{p3-83-c31-s1:eq:vacuum-reconstruct-log}
\end{align}
En particular, la orientaci\'on se pierde si se conserva s\'olo el producto,
pero se recupera al a\~nadir el cociente.
\end{corollary}

\begin{proof}
La primera l\'inea es la ra\'iz positiva de las identidades
\cref{p3-83-c31-s1:eq:vacuum-four,p3-83-c31-s1:eq:vacuum-relations}. La segunda invierte el sistema
lineal
\(\mathcal E=\log r_-+\log r_+\),
\(\mathcal B=\log r_--\log r_+\).
\end{proof}

Este corolario expresa una propiedad central de HMT: producto y cociente son
funcionales complementarios de un mismo estado. El producto determina la
propagaci\'on com\'un, mientras que el cociente conserva la orientaci\'on
relativa de las hojas interior y exterior.

En el certificado V2,

\begin{align*}
r_+&=0.9999996285482493631786952281\ldots,\\
r_-&=0.9997241371895183641315165853\ldots,\\
\widehat\varepsilon&=0.9999992570966367027604416155\ldots,\\
\widehat\mu&=0.9994483504793269350899815559\ldots,\\
\widehat Z&=0.9997245085389372154555543466\ldots,\\
\widehat c&=1.0002763104861575261063862689\ldots.
\end{align*}

Estos numerales son el reconocimiento terminal de \cref{p3-83-c31-s1:eq:vacuum-operator}; no definen sus exponentes \(90,120\).

\section{Factorizaci\'on arm\'onica de los sectores \texorpdfstring{\(90\)}{90} y \texorpdfstring{\(120\)}{120}}

El semigrupo arm\'onico es

\[
\langle90,120\rangle=30\langle3,4\rangle.
\]

Si \(s=q^{30}\), entonces

\begin{equation}
\frac{1-q^{90}}{1-q^{120}}
=\frac{1-s^3}{1-s^4}
=\frac{1+s+s^2}{1+s+s^2+s^3}.
\label{p3-83-c31-s1:eq:three-four-completion}
\end{equation}

La identidad tambi\'en separa la respuesta de su defecto. Si
\[
F_{3\to4}(s)=\frac{1+s+s^2}{1+s+s^2+s^3},
\qquad 0<s<1,
\]
entonces
\begin{equation}
d_{3\to4}(s):=1-F_{3\to4}(s)
=\frac{s^3}{1+s+s^2+s^3}>0.
\label{p3-83-c31-s1:eq:vacuum-defect34}
\end{equation}
El t\'ermino adicional se conserva exactamente como defecto
positivo. Adem\'as,
\begin{equation}
F_{3\to4}'(s)
=-\frac{s^2(3+2s+s^2)}{(1+s+s^2+s^3)^2}<0.
\label{p3-83-c31-s1:eq:vacuum-monotonic34}
\end{equation}
Por consiguiente, el orden de los canales angulares determina un\'ivocamente
el orden de las respuestas constitutivas. La asignaci\'on de
\(\widehat\mu\) y \(\widehat\varepsilon\) a las hojas respectivas se obtiene
anal\'iticamente, sin recurrir a sus expansiones decimales.

Cada hoja del vac\'io constituye la completaci\'on exacta de tres a cuatro
t\'erminos de un mismo modo elemental \(30\). En consecuencia, permeabilidad
y permitividad se interpretan como dos funcionales cuadr\'aticos de una misma
completaci\'on evaluada sobre hojas conjugadas, y no como par\'ametros
constitutivos independientes.

Como la funci\'on \((1-q^{90})/(1-q^{120})\) decrece en \(0<q<1\) y \(q_+<q_-\), se obtiene

\[
0<r_-<r_+<1,\qquad
0<\widehat\mu<\widehat\varepsilon<1,\qquad
0<\widehat Z<1<\widehat c.
\]

Los l\'imites del funcional determinan su geometr\'ia. Cuando \(q\to0^+\),
\(F(q)\to1\): las cuatro lecturas se acercan a la unidad y el defecto de
memoria desaparece. Cuando \(q\to1^-\),
\[
F(q)\longrightarrow\frac{90}{120}=\frac34.
\]
La respuesta toma, por tanto, valores en el intervalo positivo \((3/4,1)\). El
extremo \(3/4\) no es una constante ajustada: es el cociente de las dos
longitudes de canal y la realizaci\'on continua de la completaci\'on
tres--a--cuatro.

\section{Involuci\'on de hojas y conservaci\'on de la orientaci\'on}

El intercambio de hojas \(C^*\mapsto-C^*\) act\'ua por

\begin{equation}
\widehat\mu\longleftrightarrow\widehat\varepsilon,\qquad
\widehat Z\longmapsto\widehat Z^{-1},
\qquad
\widehat c\longmapsto\widehat c.
\label{p3-83-c31-s1:eq:vacuum-sheet}
\end{equation}

La propagaci\'on com\'un es par; la impedancia conserva la orientaci\'on.
Esta distinci\'on reaparecer\'a en Barbero y CKM: no todos los funcionales
asociados al cambio de hoja tienen la misma paridad.

\section{Caracteres logar\'itmicos y compatibilidad tracial}

Definamos

\[
\mathcal E=\log(r_+r_-),\qquad
\mathcal B=\log(r_-/r_+).
\]

Con la traza normalizada \(\tau=\Tr/2\),

\begin{equation}
2\tau(\log\mathcal V)=\mathcal E,
\qquad
-2\tau(R\log\mathcal V)=\mathcal B.
\label{p3-83-c31-s1:eq:vacuum-traces}
\end{equation}

La pareja \((\mathcal E,\mathcal B)\) constituye un sistema de coordenadas lineal del logaritmo
operatorio:
\begin{equation}
\log\mathcal V
=\frac{\mathcal E}{2}I-\frac{\mathcal B}{2}R.
\label{p3-83-c31-s1:eq:vacuum-log-decomposition}
\end{equation}
La componente escalar es invariante bajo el intercambio de hojas; la
componente orientada cambia de signo. Esta descomposici\'on demuestra
directamente la paridad de la propagaci\'on y el car\'acter orientado de la
impedancia.

La amplificaci\'on non\'adica

\[
\mathcal V_m=\mathcal V\otimes I_{9^m}
\]

es compatible con las expectativas \(E_m=\mathrm{id}\otimes\tau_9\):

\[
E_m(\mathcal V_{m+1})=\mathcal V_m.
\]

Por tanto, \(\mathcal E\) y \(\mathcal B\) no son peculiaridades de una matriz \(2\times2\); permanecen en la torre UHF y en su cierre tracial \(II_1\).

M\'as precisamente, para todo polinomio \(p\) y despu\'es, por continuidad,
para toda funci\'on continua en el espectro,
\begin{equation}
E_m\bigl(p(\mathcal V_{m+1})\bigr)=p(\mathcal V_m).
\label{p3-83-c31-s1:eq:vacuum-functional-refinement}
\end{equation}
La compatibilidad no preserva solamente dos n\'umeros: preserva todo el
\'algebra funcional generada por la respuesta. En particular conserva sus
proyectores espectrales, sus potencias, su logaritmo y las formas cuadr\'aticas
que miden los defectos \cref{p3-83-c31-s1:eq:vacuum-defect34}.

\section{Realizaci\'on dimensional de las relaciones constitutivas}

La Mec\'anica Dimensional determina

\begin{align}
Z_{\rm vac}&=\widehat Z\,u_Su_Q^{-2},\\
c_{\rm vac}&=\widehat c\,u_C,\\
\mu_{\rm vac}&=\widehat\mu\,u_Su_C^{-1}u_Q^{-2},\\
\varepsilon_{\rm vac}&=\widehat\varepsilon\,u_S^{-1}u_C^{-1}u_Q^2.
\label{p3-83-c31-s1:eq:vacuum-sections}
\end{align}

Las identidades constitutivas se mantienen en las rectas:

\begin{equation}
\mu_{\rm vac}\varepsilon_{\rm vac}=c_{\rm vac}^{-2},
\qquad
\mu_{\rm vac}/\varepsilon_{\rm vac}=Z_{\rm vac}^2.
\end{equation}

La carta SI es posterior. No es necesario declarar \(c\) primitivo y resolver despu\'es \(\mu\varepsilon=c^{-2}\): los tres objetos descienden conjuntamente de \(\mathcal V\) y de las rectas declaradas.

La separaci\'on de niveles se expresa mediante el diagrama
\[
\begin{array}{c}
(A,C^*)\longrightarrow T\longrightarrow\mathcal V
\longrightarrow
(\widehat\mu,\widehat\varepsilon,\widehat Z,\widehat c)
\\[2mm]
\Big\downarrow\text{ cartas de recta}\hspace{36mm}
\Big\downarrow\text{ evaluaci\'on final}
\\[2mm]
(u_S,u_Q,u_C)\longrightarrow
(\mu_{\rm vac},\varepsilon_{\rm vac},Z_{\rm vac},c_{\rm vac})
\longrightarrow\text{ coordenadas SI}.
\end{array}
\]
La primera l\'inea es adimensional y operatoria; la segunda determina tipos de
magnitud; la tercera selecciona una carta metrol\'ogica. Alterar la carta final
no modifica ni los autovalores ni las identidades constitutivas. En ese
sentido preciso, la unificaci\'on matem\'atica--f\'isica consiste en distinguir
el invariante geneal\'ogico de su sistema de coordenadas, y no en suprimir la
estructura dimensional.

\begin{proposition}[Minimalidad espectral]
Sea \(W\) un operador positivo de dos hojas con espectro simple
\(\{r_+,r_-\}\). Si cuatro caracteres positivos satisfacen
\(\mu\varepsilon=c^{-2}\) y \(\mu/\varepsilon=Z^2\), si la propagaci\'on
com\'un se determina mediante \(c^{-1}=\det W=r_-r_+\), y si su orientaci\'on
fija \(Z=r_-/r_+\), entonces necesariamente
\[
(\mu,\varepsilon,Z,c)
=\left(r_-^2,r_+^2,\frac{r_-}{r_+},\frac1{r_-r_+}\right).
\]
Por tanto \cref{p3-83-c31-s1:eq:vacuum-four} no es una parametrizaci\'on entre muchas:
es la realizaci\'on positiva m\'inima compatible con producto, cociente y
orientaci\'on.
\end{proposition}

\begin{proof}
De \(Z=r_-/r_+\) y \(c^{-1}=\det W=r_-r_+\) se obtienen de forma \`unica las
ra\'ices positivas. Las dos identidades constitutivas fuerzan despu\'es
\(\mu=r_-^2\) y \(\varepsilon=r_+^2\).
\end{proof}


\subsection{Operador constitutivo del vacío y transducción dimensional}

Sea \(R=R^*=R^{-1}\) la involución de hojas y
\[
P_\pm=\frac{I\pm R}{2},
\qquad
q_\pm=\exp\!\left[-\frac{\pi}{180}(A\pm\Cstar)\right],
\qquad T=q_+P_++q_-P_-.
\]
La respuesta bicapa es
\begin{equation}
\mathcal V_{90,120}=(I-T^{90})(I-T^{120})^{-1}
=r_+P_++r_-P_-,
\qquad
r_\pm=\frac{1-q_\pm^{90}}{1-q_\pm^{120}}.
\label{p3-83-c31-s2:eq:vacuum-operator}
\end{equation}
Las cuatro respuestas adimensionales son funciones del mismo operador:
\begin{equation}
\boxed{
\widehat\varepsilon=r_+^2,
\qquad \widehat\mu=r_-^2,
\qquad \widehat Z=\frac{r_-}{r_+},
\qquad \widehat c=\frac1{r_+r_-}.}
\label{p3-83-c31-s2:eq:vacuum-four}
\end{equation}
Satisfacen exactamente
\begin{equation}
\widehat\mu\,\widehat\varepsilon=\widehat c^{-2},
\qquad
\widehat\mu/\widehat\varepsilon=\widehat Z^2,
\label{p3-83-c31-s2:eq:vacuum-identities}
\end{equation}
y toman los valores
\[
\begin{aligned}
\widehat\varepsilon&=0.9999992570966367027604416155\ldots,\\
\widehat\mu&=0.9994483504793269350899815559\ldots,\\
\widehat Z&=0.9997245085389372154555543466\ldots,\\
\widehat c&=1.0002763104861575261063862689\ldots .
\end{aligned}
\]

La transducción hacia magnitudes dimensionales es el mapa
\begin{equation}
\mathfrak D_{u_S,u_Q,u_C}:
(\widehat\varepsilon,\widehat\mu,\widehat Z,\widehat c)
\longmapsto
(\varepsilon_{\rm vac},\mu_{\rm vac},Z_{\rm vac},c_{\rm vac}),
\label{p3-83-c31-s2:eq:vacuum-transducer}
\end{equation}
definido por
\begin{equation}
\boxed{
\begin{aligned}
\varepsilon_{\rm vac}&=\widehat\varepsilon\,u_S^{-1}u_C^{-1}u_Q^2,\\
\mu_{\rm vac}&=\widehat\mu\,u_Su_C^{-1}u_Q^{-2},\\
Z_{\rm vac}&=\widehat Z\,u_Su_Q^{-2},\\
c_{\rm vac}&=\widehat c\,u_C.
\end{aligned}}
\label{p3-83-c31-s2:eq:vacuum-dimensional-sections}
\end{equation}
Por homogeneidad,
\[
\mu_{\rm vac}\varepsilon_{\rm vac}=c_{\rm vac}^{-2},
\qquad
\mu_{\rm vac}/\varepsilon_{\rm vac}=Z_{\rm vac}^2.
\]
Sólo después actúa la carta \(\chi_{\rm SI}\). De este modo no se iguala
ningún carácter adimensional con una magnitud SI.

\begin{center}
\small
\begin{tabularx}{.98\textwidth}{@{}p{.13\textwidth}Y Y@{}}
\toprule
Objeto & Sección dimensional & Coordenada en la carta \(\chi_{\rm SI}\)\\
\midrule
Permitividad & \(\varepsilon_{\rm vac}=\widehat\varepsilon u_S^{-1}u_C^{-1}u_Q^2\) &
\(\varepsilon_0=8.854187812793002329\ldots\times10^{-12}\,\mathrm{F\,m^{-1}}\)\\
Permeabilidad & \(\mu_{\rm vac}=\widehat\mu u_Su_C^{-1}u_Q^{-2}\) &
\(\mu_0=1.256637062121047789\ldots\times10^{-6}\,\mathrm{H\,m^{-1}}\)\\
Impedancia & \(Z_{\rm vac}=\widehat Z u_Su_Q^{-2}\) &
\(Z_0=376.730313667167610257\ldots\,\Omega\)\\
Propagación & \(c_{\rm vac}=\widehat c u_C\) &
\(c_0=299792458\,\mathrm{m\,s^{-1}}\)\\
\bottomrule
\end{tabularx}
\end{center}

En la carta SI posterior a 2019, las coordenadas definitorias de \(e,h,c\) y
la sección \(\alphaHMT\) publican el mismo vacío mediante
\begin{equation}
\boxed{
Z_0=\frac{2h\alphaHMT}{e^2},\qquad
\mu_0=\frac{Z_0}{c},\qquad
\varepsilon_0=\frac1{Z_0c},\qquad
(\mu_0\varepsilon_0)^{-1/2}=c.}
\label{p3-83-c31-s2:eq:vacuum-si-closure}
\end{equation}
No hay dos juegos rivales de constantes: los caracteres con acento
circunflejo describen la respuesta espectral interna, las secciones con
subíndice \({\rm vac}\) conservan su tipo dimensional y los símbolos con
subíndice \(0\) son sus coordenadas en la carta física.

\subsection{Tomografía reversible del vacío}

La respuesta constitutiva no sólo desciende del operador angular: permite
reconstruirlo. Definamos
\begin{equation}
F(q)=\frac{1-q^{90}}{1-q^{120}},\qquad 0<q<1.
\label{p3-83-c31-s2:eq:vacuum-F}
\end{equation}
Con \(s=q^{30}\),
\[
F(q)=\frac{1+s+s^2}{1+s+s^2+s^3},
\qquad
\frac{\dd F}{\dd s}
=-\frac{s^2(3+2s+s^2)}{(1+s+s^2+s^3)^2}<0.
\]
Los límites \(F(0^+)=1\) y \(F(1^-)=3/4\) demuestran el difeomorfismo
\begin{equation}
\boxed{F:(0,1)\xrightarrow{\;\sim\;}(3/4,1).}
\label{p3-83-c31-s2:eq:vacuum-F-bijection}
\end{equation}
Para cada \(r\in(3/4,1)\), la inversa es
\(F^{-1}(r)=s^{1/30}\), donde \(s\in(0,1)\) es la raíz única de
\begin{equation}
rs^3+(r-1)(1+s+s^2)=0.
\label{p3-83-c31-s2:eq:vacuum-inverse-cubic}
\end{equation}

La inversión simultánea de las dos hojas recupera el antecedente completo:
\begin{equation}
\boxed{
\begin{gathered}
q_+=F^{-1}(\sqrt{\widehat\varepsilon}),\qquad
q_-=F^{-1}(\sqrt{\widehat\mu}),\qquad
D_A=q_+q_-,\\
A_{\rm rad}=-\frac12\log D_A,\qquad
C^*_{\rm rad}=\frac12\log\frac{q_-}{q_+},
\qquad T=F^{-1}(\mathcal V_{90,120}).
\end{gathered}}
\label{p3-83-c31-s2:eq:vacuum-tomography}
\end{equation}
Vacío y par angular son dos coordenadas reversibles del mismo objeto. El
producto de hojas publica la propagación común; el cociente conserva su
orientación. Esta tomografía elimina la ambigüedad entre ``valor interno'' y
``valor físico'': \(\widehat\varepsilon,\widehat\mu\) son caracteres del
operador, \(\varepsilon_{\rm vac},\mu_{\rm vac}\) son sus secciones
dimensionales y \(\varepsilon_0,\mu_0\) son sus coordenadas en la carta SI.

\subsubsection{Caracteres logarítmicos y estabilidad nonádica}

La respuesta bicapa separa su componente común de la orientación de hoja.
Para la traza normalizada \(\tau_2\) de \(M_2(\mathbb C)\),
\begin{equation}
\boxed{
\mathcal E=2\tau_2(\log\mathcal V_{90,120}),\qquad
\mathcal B=-2\tau_2(R\log\mathcal V_{90,120}),\qquad
\log\mathcal V_{90,120}
=\frac{\mathcal E}{2}I-\frac{\mathcal B}{2}R.}
\label{p3-83-c31-s2:eq:vacuum-log-characters}
\end{equation}
El carácter \(\mathcal E\) publica la respuesta común y \(\mathcal B\)
conserva su paridad. Esta descomposición no depende del tamaño
\(2\times2\). En el refinamiento nonádico, sean
\begin{equation}
\begin{gathered}
\mathfrak A_m=M_2(\mathbb C)\otimes M_{9^m}(\mathbb C),\qquad
\mathfrak A_{m+1}\cong\mathfrak A_m\otimes M_9(\mathbb C),\\
\iota_m(x)=x\otimes I_9,\qquad
E_m=\operatorname{id}_{\mathfrak A_m}\otimes\tau_9:
\mathfrak A_{m+1}\longrightarrow\mathfrak A_m,\\
\mathcal V_m=\mathcal V_{90,120}\otimes I_{9^m},\qquad
E_m(\mathcal V_{m+1})=\mathcal V_m.
\end{gathered}
\label{p3-83-c31-s2:eq:vacuum-uhf-tower}
\end{equation}
Aquí \(\tau_9\) es la traza normalizada de \(M_9(\mathbb C)\). Con
\(\tau_m=\tau_2\otimes\tau_{9^m}\) y
\(R_m=R\otimes I_{9^m}\), los caracteres sobreviven literalmente:
\begin{equation}
\mathcal E_m=2\tau_m(\log\mathcal V_m)=\mathcal E,
\qquad
\mathcal B_m=-2\tau_m(R_m\log\mathcal V_m)=\mathcal B.
\label{p3-83-c31-s2:eq:vacuum-uhf-characters}
\end{equation}
Además, para todo polinomio \(p\),
\begin{equation}
E_m\!\left(p(\mathcal V_{m+1})\right)=p(\mathcal V_m).
\label{p3-83-c31-s2:eq:vacuum-conditional-polynomial}
\end{equation}
El cálculo funcional continuo prolonga la compatibilidad a proyectores,
potencias y logaritmos. Operador y caracteres sobreviven así en el límite
UHF y en su cierre factorial tracial de tipo \(II_1\) como invariantes de la
torre completa.
